\documentclass{article}
\usepackage[margin=0.75in, top=0.6in, bottom=0.6in]{geometry}
\usepackage{enumitem}

\title{\vspace{-1.5em}Final Project Proposal\\[0.3em]\large Procedural Geographic Map Generation via Generative Adversarial Networks\vspace{-1em}}
\author{}
\date{}

\begin{document}
\maketitle

\section*{Summary}
We will train a GAN to produce terrain maps for use in a simple video game. 
The training maps will be generated from Perlin-noise, so the GAN will be learning
the Perlin-noise generation process.

\section*{Description}

\paragraph{Approach.}
A DCGAN maps a latent vector $\mathbf{z} \sim \mathcal{N}(0,I)$ to a fixed-resolution bitmap
where each pixel encodes a terrain type (water, grassland, forest, mountain, etc.).

\paragraph{Data.}
The training set consists of \textbf{2,000--5,000} maps generated via layered simplex (Perlin) noise
with terrain-type thresholds applied to assign each pixel one of several categories
(water, grassland, forest, mountain, etc.). The GAN's objective is to learn the distribution
induced by this noise process, which enables the generation of novel terrain maps.

\paragraph{Evaluation.}
Image quality will be measured with Fr\'echet Inception Distance (FID).
Since the training distribution is Perlin noise, frequency-domain statistics (power spectrum slope)
of generated maps can be compared directly against training samples as a distribution-level check.
Navigability will be assessed with two geometry-based metrics derived from A* traversal:
\textbf{path tortuosity} (ratio of shortest-path length to Euclidean distance between endpoints,
averaged over random pairs) measures how geographically interesting the terrain is, and
\textbf{reachable area fraction} (fraction of traversable pixels reachable from a random start)
measures connectivity and fragmentation.
Together these reward maps that are neither trivially open nor hopelessly disconnected.

\section*{Project Components}

\begin{itemize}[leftmargin=1.5em, topsep=2pt, itemsep=0pt]
    \item \textbf{True Data generation:} generate maps via layered Perlin noise; different noise levels produce different colors representing terrain types.
    \item \textbf{Baseline DCGAN:} generator and discriminator; verify stability with loss curves and sample grids.
    \item \textbf{Navigability metrics:} A* traversal; compute path tortuosity and reachable area fraction on generated maps.
    \item \textbf{Quantitative evaluation:} FID, power spectrum comparison, and navigability metrics on a held-out test set.
    \item \textbf{Game development and integration:} develop a very simple navigation game, with live, in-game map generation.
    \item \textbf{Written report:} method, experiments, results, and discussion.
    \item \textbf{Presentation / demo:} live or recorded in-game demo, slideshow summarizing approach and results.
\end{itemize}

\vspace{1em}
\noindent\small\textit{Assistance from Claude AI.}

\end{document}
